% Generated draft: requires mathematical and source review.
\documentclass[11pt]{article}
\usepackage[margin=25mm]{geometry}
\usepackage{fontspec}
\setmainfont{Latin Modern Roman}
\usepackage{amsmath,amssymb,mathtools,unicode-math}
\setmathfont{Latin Modern Math}
\usepackage{enumitem,needspace,xurl,hyperref}
\setlength{\parindent}{0pt}
\setlength{\parskip}{5pt}
\setlength{\emergencystretch}{3em}
\newcommand{\N}{\mathbb{N}}
\newcommand{\Z}{\mathbb{Z}}
\newcommand{\Q}{\mathbb{Q}}
\newcommand{\R}{\mathbb{R}}
\newcommand{\C}{\mathbb{C}}
\newcommand{\F}{\mathbb{F}}
\newcommand{\A}{\mathbb{A}}
\newcommand{\PP}{\mathbb{P}}
\newcommand{\E}{\mathbb{E}}
\newcommand{\eps}{\varepsilon}
\newcommand{\dd}{\,\mathrm{d}}
\begin{document}
\section*{UnsolvedMath ID: 30002704}
\textbf{Source problem code:} catalogue entry
\subsection{Definitions and mathematical statement}
Write \(\mathcal M(X)\) for the source's space of Radon measures on \(X\), and let \(G\) be a group acting on this space. The source does not specify the topology of \(X\), the precise convention for Radon measures, or the measurable structure on \(\mathcal M(X)\). Consequently, the following definition is conditional on supplying a sigma-algebra \(\Sigma\) on \(\mathcal M(X)\) for which each action map \(T_g:\mu\mapsto g\cdot\mu\), \(g\in G\), is measurable. No particular choice of these missing structures is asserted here.

A law of a random measure is a probability measure \(P\) on \((\mathcal M(X),\Sigma)\). Its pushforward under \(g\in G\) is defined by
\[
 (T_g)_*P(A)=P(T_g^{-1}(A)),\qquad A\in\Sigma.
\]
The law \(P\) is \emph{quasi-invariant under \(G\)} if
\[
 \forall g\in G,\qquad (T_g)_*P\sim P,
\]
where equivalence means mutual absolute continuity; explicitly,
\[
 \forall g\in G\ \forall A\in\Sigma,\qquad
 P(A)=0\ \Longleftrightarrow\ P(T_g^{-1}(A))=0.
\]

\textbf{Source question.} For a group \(G\) acting on \(\mathcal M(X)\), which probability laws \(P\) are quasi-invariant under \(G\) and support an associated analysis, such as a Laplace operator or a diffusion on \(\mathcal M(X)\)?

The source poses a general research question rather than a fully specified classification problem. It supplies neither a class of spaces and actions to be classified nor a definition of the required associated analysis. In particular, it does not specify the domain or properties of a Laplace operator, the topology and path properties of a diffusion, or how either object must be related to \(P\) and the action. The examples of a Laplace operator and a diffusion are illustrative; the source does not require both. These missing qualifications prevent a unique, fully precise mathematical target from being extracted.
\subsection{Short English statement}
The question asks which random-measure laws preserve their null sets under every element of a given group action and also admit useful associated analysis, for example a Laplace operator or a diffusion. Quasi-invariance has a precise conditional formulation above, but the underlying spaces, actions, and analytic requirements remain unspecified in the source.
\subsection{Sources}
\begin{itemize}
\item Supplied local UnsolvedMath registry record, permanent ID \(30002704\), titled \emph{Quasi-Invariant Random Measures and Associated Diffusions}; exact source problem code \texttt{catalogue entry}. Collection: \emph{Oberwolfach Reports Open Problems}. The record gives proposed year \(2014\) and identifies no individual proposer.
\item Supplied report citation: \emph{Dirichlet Form Theory and its Applications (2015)}, Oberwolfach Reports, pp.~2667--2756, DOI \texttt{10.4171/owr/2014/48}. Supplied source URL: \url{https://doi.org/10.4171/owr/2014/48}. The background describes Oberwolfach Reports as documenting mathematical workshops and their open-problem discussions; it supplies no narrower problem location within the report.
\item Supplied literature source: \emph{Oberwolfach Report 48/2014}, \url{https://doi.org/10.4171/owr/2014/48}. The exact supplied literature-check date is \texttt{2026-08-21}. This is supplied dated curation, not independently verified: the assessment labels the broad question open and describes a request for a general theory without a precise class or theorem. The background explicitly calls for independent verification of time-sensitive status and cited claims.
\item The registry labels the statement \texttt{corrected\_verified}, with a supplied claim of manual semantic, typographic, and LaTeX review of a meaning-preserving cleanup. This is recorded provenance, not independent verification of the report attribution, mathematical formulation, or open status.
\end{itemize}
\end{document}
